\chapter{The matching engine at general arity}
\label{ch:engine}

The matching theorem for tree-indexed Markov label fields whose counter laws share a common
core, with the one-law theorem as a corollary. Formalised in the library
\texttt{GraphMarkovMatching}, self-contained over its vendored support layer.

\section{The matching theorems}

\begin{theorem}[Two-law matching]
\label{thm:composite-matching}
\lean{GraphMarkovMatching.Composite.composite_failure_bound_massfree}
\leanok
\uses{def:etaG}
Let two counter laws share a common core. Under the potential hypothesis on the label graph,
the probability that no automorphism matches the two tree-indexed label fields is bounded by an
explicit multiple of the potential, uniformly in the height, and the same bound holds along the
infinite tree.
\end{theorem}

\begin{corollary}[One-law matching]
\label{thm:main-matching}
\lean{GraphMarkovMatching.varying_failure_le}
\leanok
Specialising both sides to the same counter law gives the one-law matching theorem, with the
failure probability bounded by $K_c\eta$ for an explicit constant.
\end{corollary}

\begin{proof}
\leanok
\uses{thm:composite-matching, thm:numeric}
Specialise the two counter laws of \cref{thm:composite-matching} to a single law. The tagged
grammar collapses to the ordinary one, the marker letters never occur, and the constants of
\cref{thm:numeric} give the bound $K_c\eta$.
\end{proof}

\begin{corollary}[Exponential three-label example]
\label{thm:exponential}
\lean{GraphMarkovMatching.varying_exp3_failure_le}
\leanok
Fix $N\geq2$ and a law $\nu$ supported on a nonempty subset of $\set{2,\dots,N}$, let $G$ be the
path $0-1-2$, and for real $D\geq5$ let $\mu_D$ have masses $p_1=e^{-D}$, $p_2=e^{-D^2}$ and
$p_0=1-p_1-p_2$. Then the matching hypothesis holds for all large $D$.
\end{corollary}

\begin{proof}
\leanok
\uses{thm:main-matching}
Every state is within distance one of $1$, with $b(0)=1-p_2$, $b(1)=1$ and $b(2)=p_1+p_2$. For
$D\geq5$ both masses are below $e^{-5}<\tfrac1{20}$, so the tilts are bounded by a constant and the
potential by a multiple of $e^{-D}$, which meets the hypothesis for large $D$.
\end{proof}

\begin{corollary}[Double-exponential law on the half-line]
\label{thm:doubleexp}
\lean{GraphMarkovMatching.varying_doubleExp_failure_le}
\leanok
The same holds for the path on the nonnegative integers with $p_k=e^{-D^k}$ for $k\geq1$ and
$p_0=1-\sum_{k\geq1}p_k$.
\end{corollary}

\begin{proof}
\leanok
\uses{thm:main-matching, thm:double-exp}
For $k\geq1$ the ratio $p_{k+1}/p_k=e^{-D^k(D-1)}\leq e^{-20}$ makes every tail geometric, so
$\sum_{k\geq1}p_k<\tfrac1{10}$ and $p_0\geq\tfrac12$. Local dominance on the half-line then bounds
the potential by a multiple of $e^{-D(D-5/2)}$.
\end{proof}

\section{Potential-theoretic estimates}

\begin{lemma}[Zero interface]
\label{thm:zero-interface}
\lean{GraphMarkovMatching.qE_zero_le}
\leanok
\uses{def:etaG}
Suppose $\mu(0)\geq1/2$ and write $\eta:=\eta_{G,\alpha}(\mu)$. Then
$\mu(V\setminus B_G(0,1))\leq2\eta$, and the far tilt and tilt moment sums are bounded by
explicit multiples of $\eta$.
\end{lemma}

\begin{proof}
\leanok
Assume $\eta<\infty$, so every state of positive mass has $b(v)>0$. Each state outside
$B_G(0,1)$ has $b(v)\leq1-\mu(0)\leq\tfrac12$, so its summand in $\eta$ is at least
$\tfrac12\mu(v)$, which gives the mass bound. The tilt and moment sums follow by the same pointwise
comparison against the summands of $\eta$.
\end{proof}

\begin{theorem}[Four-law inequality]
\label{thm:four-law}
\lean{GraphMarkovMatching.PhiD_fourlaw_le}
\leanok
\uses{def:potential, def:square}
Let $\alpha\geq1$ and $\delta,L,K>0$ satisfy the standing hypothesis, let $R$ be a symmetric
relation on a countable set, and let $\mu_0,\mu_1,\nu_0,\nu_1$ be laws. If $M$ bounds all eight
directed potentials between the two source and two target laws, then the potential of the
symmetrised square of the pair laws is bounded by an explicit function of $M$.
\end{theorem}

\begin{proof}
\leanok
\uses{thm:pointwise-fourlaw}
Sort the two rows in each column and apply the pointwise bound of \cref{thm:pointwise-fourlaw}
at every doubly-positive pair. Integrating against the product of the two source laws and bounding
each of the eight directed potentials by $M$ collects the linear and quadratic terms into the stated
constants.
\end{proof}

\begin{proposition}[Pointwise four-law bound]
\label{thm:pointwise-fourlaw}
\lean{GraphMarkovMatching.phi_Q_split_four}
\leanok
At a fixed pair $(x_0,x_1)$ with all four bad degrees below one, the square bad degree is below
one and the weight of the square is bounded by a half-sum of the four one-sided weights plus a
quadratic remainder.
\end{proposition}

\begin{proof}
\leanok
\uses{def:potential}
In each column put $d_j=\max(q_0^j,q_1^j)$ and $e_j=\min(q_0^j,q_1^j)$. A lower bound for the
square good degree in terms of the four column quantities, followed by the tangent-line and chord
estimates for $\phi^E_\alpha$, separates the half-sum of the four one-sided weights from a quadratic
remainder.
\end{proof}

\section{The one-law transfer system}

\begin{proposition}[The frozen cross-type potential is infinite]
\label{thm:obstruction}
\lean{GraphMarkovMatching.crossed_context_potential_top}
\leanok
\uses{def:potential}
Let $\nu=\delta_3$ and let $v_f$ be any state with $\mu(v_f)>0$ and $v_f\not\sim0$. Then the
directed potential from the frozen cross type into the fresh law is infinite at every height.
\end{proposition}

\begin{proof}
\leanok
Build, for each height $m$, an incompatible-root labelling together with an auxiliary labelling,
by grafting the frozen state $v_f$ at the root and descending along the $\delta_3$ cascade. The
source mass stays positive while the target good degree vanishes at every such point, so the
directed potential has an infinite summand at every height.
\end{proof}

This is why the recursion is carried over screens rather than over the naive cross types.

\begin{lemma}[One-step decomposition]
\label{thm:one-step}
\lean{GraphMarkovMatching.rE_succ_branch}
\leanok
Fix a height-$(h+1)$ source point $x=\mathrm{branch}((v,k);x^1,x^2)$ with root state $v$ and
child pair $xp$. The good and bad degrees at $x$ decompose over the root state and the two
children.
\end{lemma}

\begin{proof}
\leanok
The target law at height $h+1$ is the pushforward of the child-pair law under the branch map.
Computing the good degree through that pushforward and applying the branch recursion splits it into
the root-state factor and the two child factors.
\end{proof}

\begin{lemma}[Screening a fresh mixture]
\label{thm:fresh-mixture}
\lean{GraphMarkovMatching.rE_bind}
\leanok
Let $\rho=\sum_a w_a\rho_a$ be a countable mixture of laws. Then $r_\rho(z)=\sum_a w_ar_{\rho_a}(z)$
for every point $z$, so $r_\rho(z)=0$ if and only if $r_{\rho_a}(z)=0$ for every component with
$w_a>0$.
\end{lemma}

\begin{proof}
\leanok
The good degree is a sum of target masses over compatible points, hence linear in the target
law, and Fubini for nonnegative sums gives the mixture identity. A vanishing sum of nonnegative
terms forces each term with positive weight to vanish.
\end{proof}

\begin{definition}[Restricted quantities]
\label{def:restricted}
\lean{GraphMarkovMatching.PhiDres}
\leanok
\uses{def:potential}
For laws $\rho_s,\rho_t$ and a relation $R$, the restricted potential sums the weight only over
points at which the target good degree is positive; the restricted weight and the zero mass are
defined on the same restriction.
\end{definition}

\begin{definition}[Normalised screens]
\label{def:screen}
\lean{GraphMarkovMatching.screenE}
\leanok
\uses{def:restricted}
Let $\rho_s$ be a cell law, $\boldsymbol z$ a finite list of target laws, and
$g\colon\cX\to[0,\infty]$ a normalisation. The screen $\cE(\rho_s;\boldsymbol z\!\not\to g)$
integrates $g$ against $\rho_s$ over the points that are dead against every member of
$\boldsymbol z$.
\end{definition}

\begin{lemma}[Diagonal pruning]
\label{thm:pruning}
\lean{GraphMarkovMatching.screenE_self_mem}
\leanok
\uses{def:screen}
If the cell law is itself a member of the zero list and $R$ is reflexive, then the screen
vanishes for every normalisation.
\end{lemma}

\begin{proof}
\leanok
Fix $x$ with $\rho_s(x)>0$. Reflexivity gives $r_{\rho_s}(x)\geq\rho_s(x)>0$, and since
$\rho_s$ is in the zero list the dead-event indicator vanishes at $x$. Every summand of the screen
is therefore zero.
\end{proof}

\begin{lemma}[Tilt pruning]
\label{thm:tilt-pruning}
\lean{GraphMarkovMatching.screenE_tilt_self}
\leanok
\uses{def:screen}
If the normalisation is the restricted tilt of a member of the zero list, the screen vanishes.
\end{lemma}

\begin{proof}
\leanok
Fix $x$. If $r_\rho(x)=0$ the restricted weight vanishes there by definition; if $r_\rho(x)>0$
the dead-event indicator vanishes because $\rho$ lies in the zero list. Either way every summand is
zero.
\end{proof}

\begin{lemma}[Screened $2\times2$ classification]
\label{thm:hall}
\lean{GraphMarkovMatching.rE_square_eq_zero_iff_hall}
\leanok
\uses{def:screen}
The square good degree vanishes exactly when one source child is dead against both target laws,
or the Hall condition fails on the $2\times2$ array; the surviving cases split into an explicit
sum of resolved terms.
\end{lemma}

\begin{proof}
\leanok
\uses{thm:one-step}
The square compatibility event is the union of the straight and the crossed pairing events, so
its degree vanishes exactly when both pairings fail. Reading the four one-sided degrees as a
$2\times2$ array, that is the failure of the Hall condition, and the surviving cases are the
resolved terms.
\end{proof}

\begin{lemma}[Mixture bound for restricted weights]
\label{thm:mixture-tilt}
\lean{GraphMarkovMatching.WresD_XiBar_le_sum}
\leanok
\uses{def:restricted}
For a pair of height-$h$ full labellings, the restricted weight of the mixed pair law is bounded
by the tilt-weighted sum of the restricted weights of its components.
\end{lemma}

\begin{proof}
\leanok
\uses{thm:fresh-mixture}
If the mixed good degree vanishes the left side is zero. Otherwise the mixture identity of
\cref{thm:fresh-mixture} writes the mixed degree as the weighted sum of the component degrees, and
dividing through gives the tilt-weighted bound.
\end{proof}

\section{The grammar and the accessible index}

\begin{definition}[Targets, cascade pairs, successors]
\label{def:grammar}
\lean{GraphMarkovMatching.Tgt}
\leanok
The target alphabet $\cA$ consists of the forced letters $\Zf_j$, the fresh letter $\Ff$, and
the component letters $\Ff_k$. The cascade pair of a counter and the successor relation on
letters are defined by the kernel recursion.
\end{definition}

\begin{lemma}[The bounded alphabet and screen family]
\label{thm:universes}
\lean{GraphMarkovMatching.screenUniv_card_le}
\leanok
\uses{def:grammar}
Fix $N\geq2$ with $S_\nu\subseteq\set{0,\dots,N}$. Then successors stay inside $\cA_N$, full
successors of members of $\cU_N$ stay in $\cU_N$, and
$\abs{\cA_N}\leq2N+3$, $\abs{\cU_N}\leq(2N+3)2^{2N+3}(2N+4)$.
\end{lemma}

\begin{proof}
\leanok
Every member of the cascade pair of $j\leq N$ is $\Ff$ or a letter $\Zf_i$ with
$i\leq\max\set{\floor{j/2},\,j-\floor{j/2}}\leq N$, so successors stay inside $\cA_N$. Counting the
three families of letters gives the alphabet bound, and a screen is a cell, a subset of the alphabet
and a normalisation, which gives the second count.
\end{proof}

\begin{definition}[Interpretation of the grammar]
\label{def:interp}
\lean{GraphMarkovMatching.interpT, GraphMarkovMatching.interpScreen}
\leanok
\uses{def:grammar}
At height $h$ the letters are interpreted as the process laws $Z_{j,h}$, $T_h$ and $F_{k,h}$,
and a normalisation letter as the corresponding restricted tilt.
\end{definition}

\begin{definition}[Accessible ordered pairs]
\label{def:ordacc}
\lean{GraphMarkovMatching.OrdAcc}
\leanok
\uses{def:grammar}
The set $\mathfrak O_\nu$ of accessible ordered pairs is the smallest subset of $\cA^2$
containing $(\Ff,\Ff)$ and closed under pair successors.
\end{definition}

\begin{lemma}[Invariants of accessible pairs]
\label{thm:ordacc-inv}
\lean{GraphMarkovMatching.OrdAcc.swap}
\leanok
\uses{def:ordacc}
$\mathfrak O_\nu$ is closed under swapping coordinates and under the two diagonal maps, no
component letter occurs in an accessible pair, and $\mathfrak O_\nu\subseteq\cA_N^2$ when
$S_\nu\subseteq\set{0,\dots,N}$.
\end{lemma}

\begin{proof}
\leanok
Induct on the construction. The initial pair is fixed by swapping and by both diagonal maps, and
each closure property is preserved by passing to pair successors. No successor of a letter in an
accessible pair is a component letter, so component letters never appear.
\end{proof}

\begin{definition}[Accessible screens and the live index]
\label{def:scracc}
\lean{GraphMarkovMatching.ScrAcc}
\leanok
\uses{def:screen, def:ordacc}
A formal screen is a seed when its cell, zero list and normalisation all come from the
successors of one accessible pair. The live index $\mathfrak I_\nu$ consists of the live members
of the closure of the seeds under the successor operations.
\end{definition}

\begin{lemma}[Invariants of accessible screens]
\label{thm:scracc-inv}
\lean{GraphMarkovMatching.ScrAcc.nonempty}
\leanok
\uses{def:scracc}
Every accessible screen has a nonempty zero list, no component letter occurs as its cell, in its
zero list, or as its normalisation, and its cell pairs accessibly with itself and with its
normalisation.
\end{lemma}

\begin{proof}
\leanok
\uses{thm:ordacc-inv}
Induct on the construction of a screen. A seed has a nonempty zero list by definition and each
of its letters is a successor of a coordinate of an accessible pair, so no component letter occurs;
both properties are preserved by the closure operations. The pairing claims follow from
\cref{thm:ordacc-inv}.
\end{proof}

\begin{definition}[The block matrix]
\label{def:accN}
\lean{GraphMarkovMatching.accN}
\leanok
\uses{def:scracc}
Fix $C_W\in[0,\infty]$. The block matrix on the live index carries $C_W$ from $i$ to each full
successor $j$ of $i$, and zero elsewhere.
\end{definition}

\section{The transfer rows}

\begin{lemma}[Restricted four-law]
\label{thm:res-four-law}
\lean{GraphMarkovMatching.PhiDres_fourlaw_le}
\leanok
\uses{def:restricted}
The four-law inequality holds for the restricted potentials, with an added zero-mass term.
\end{lemma}

\begin{proof}
\leanok
\uses{thm:four-law}
Run the proof of \cref{thm:four-law} on the restricted quantities. The pairs outside the
doubly-positive set contribute the zero-mass term, and the rest is bounded exactly as before.
\end{proof}

\begin{definition}[The cell bound]
\label{def:cellCB}
\lean{GraphMarkovMatching.cellCB}
\leanok
\uses{thm:res-four-law}
For bounds $M,Z,E'$ the cell bound $\mathrm{CB}(M,Z,E')$ collects the linear, quadratic,
zero-mass and screen contributions of one cell.
\end{definition}

\begin{lemma}[One-cell row]
\label{thm:one-cell}
\lean{GraphMarkovMatching.PhiDres_square_le}
\leanok
\uses{def:cellCB, def:screen}
If in addition the four resolved screens are bounded by $E'$, then the restricted potential of
the square is at most the cell bound.
\end{lemma}

\begin{proof}
\leanok
\uses{thm:res-four-law, def:screen}
The restricted square potential is the average of the weight over the pairs alive against the
target product. Splitting that average by \cref{thm:res-four-law} and bounding the four resolved
screens by $E'$ collects the terms into the cell bound.
\end{proof}

\begin{lemma}[Square cell between cascades]
\label{thm:psirow-square}
\lean{GraphMarkovMatching.psiRow_square}
\leanok
\uses{def:interp}
Fix counters $k,j$ and a height $h$. If the ordinary coordinates, zero masses and screens
between the two cascade families are bounded, then so is the square cell between them.
\end{lemma}

\begin{proof}
\leanok
\uses{thm:one-cell}
Both cells are products of interpreted letters, so \cref{thm:one-cell} applies with the two
cascade laws as source and target. The hypotheses supply the bounds it needs on the ordinary
coordinates, the zero masses and the screens.
\end{proof}

\begin{definition}[The one-sided bound]
\label{def:oneSided}
\lean{GraphMarkovMatching.oneSidedBound}
\leanok
For bounds $M,E'$ and a tilt bound $T_\nu$, the one-sided bound is
$\mathrm{OS}(M,E'):=T_\nu\abs{S_\nu}\bigl(4E'(1+\alpha M)+4E'^2\bigr)$.
\end{definition}

\begin{lemma}[One-sided tilt mass]
\label{thm:one-sided}
\lean{GraphMarkovMatching.oneSided_charge_le}
\leanok
\uses{def:oneSided}
Under the row hypotheses the one-sided charge at a counter is at most $\mathrm{OS}(M,E')$.
\end{lemma}

\begin{proof}
\leanok
\uses{thm:mixture-tilt}
Treat the indicator and the tilt separately. Some supported component is dead at a pair point
exactly when one of the finitely many component events occurs, and the tilt is bounded by
\cref{thm:mixture-tilt}. Multiplying the two bounds and summing over the support gives
$\mathrm{OS}(M,E')$.
\end{proof}

\begin{lemma}[The four ordinary rows]
\label{thm:psi-rows}
\lean{GraphMarkovMatching.psiRow_ZZ}
\leanok
\uses{def:cellCB}
With $\Delta:=\mathrm{CB}(M,Z,E')+\mathrm{OS}(M,E')$, each of the four ordinary coordinates at
height $h+1$ is bounded by $\Delta$ under the corresponding hypotheses at height $h$.
\end{lemma}

\begin{proof}
\leanok
\uses{thm:psirow-square, thm:one-sided}
Each of the four rows has a forced or fresh source whose height-$(h+1)$ law is the pushforward
of a cascade law below a compatible root. \Cref{thm:psirow-square} bounds the square cell and
\cref{thm:one-sided} the tilt mass carried by dead components, and the two together give
$\Delta$.
\end{proof}

\begin{lemma}[Descent of a zero list]
\label{thm:descent}
\lean{GraphMarkovMatching.memInd_branch_iff}
\leanok
\uses{def:screen}
For an $\Ff_k$-free zero list and a root state $v\sim0$ with $b(v)>0$, deadness of a branch
point against the list is equivalent to deadness of the child pair against the descended list.
\end{lemma}

\begin{proof}
\leanok
\uses{thm:one-step}
A member of the zero list is dead after the root decomposition exactly when its member pairs are
dead at the child pair, by the one-step decomposition and $b(v)>0$. Taking the conjunction over the
list gives the equivalence.
\end{proof}

\begin{lemma}[Multi-member Hall factorisation]
\label{thm:hall-factor}
\lean{GraphMarkovMatching.hallFactorize}
\leanok
When every component of a member of the product list lies in the scalar list, deadness against
the product list factorises through deadness against the scalar list.
\end{lemma}

\begin{proof}
\leanok
\uses{thm:hall}
Bound the indicator of joint deadness pointwise by the sum of the two one-sided indicators and a
cross term over the scalar list. Integrating that bound against the two source laws factorises the
product deadness through the scalar deadness.
\end{proof}

\begin{lemma}[The six screen rows]
\label{thm:screen-rows}
\lean{GraphMarkovMatching.interpScreen_forced_step}
\leanok
\uses{def:scracc}
For a formal screen whose cell is forced or fresh, with nonempty $\Ff_k$-free zero list and
$\Ff_k$-free normalisation, the screen at height $h+1$ is bounded by the sum of its successor
screens at height $h$ plus the ledger terms.
\end{lemma}

\begin{proof}
\leanok
\uses{thm:descent, thm:hall-factor}
For a forced cell the height-$(h+1)$ law is the pushforward of a cascade law below a compatible
root, so \cref{thm:descent} rewrites the zero event at the child pair and \cref{thm:hall-factor}
splits it over the successor screens. The fresh cell is the same computation after expanding the
root mixture, and the ledger terms collect the root tilt and the far mass.
\end{proof}

\begin{lemma}[Base bounds]
\label{thm:base}
\lean{GraphMarkovMatching.interpPhi_base}
\leanok
\uses{def:interp}
Suppose $\mu(0)\geq1/2$. At height $0$ every $\Ff_k$-free ordinary coordinate is at most
$\eta+\phi^E_\alpha(q_\mu(0))$, and every admissible formal screen is bounded likewise.
\end{lemma}

\begin{proof}
\leanok
\uses{thm:zero-interface}
A height-$0$ labelling is a single label and matching is state compatibility. For a forced
target the degree at a compatible root is $1$, so the coordinate is $\eta$ plus the weight at the
zero state; the screens are bounded by the same two quantities.
\end{proof}

\section{Nilpotence of the accessible screen block}

\begin{definition}[Return lengths]
\label{def:renewal}
\lean{GraphMarkovMatching.toFresh}
\leanok
The return lengths of a counter are defined by the recursion mirroring the kernel:
$\mathrm{toF}(k)=\set1$ for $k\leq2$, $\mathrm{toF}(3)=\set{1,2}$, and
$\mathrm{toF}(k)=(\mathrm{toF}(\floor{k/2})+1)\cup(\mathrm{toF}(\ceil{k/2})+1)$ for $k\geq4$.
\end{definition}

\begin{lemma}[Cofinality of the generated monoid]
\label{thm:semigroup}
\lean{GraphMarkovMatching.depths_semigroup}
\leanok
\uses{def:renewal}
There is a threshold $R_0$ such that every multiple of $d_\nu$ at least $R_0$ belongs to
$\Sigma_\nu$.
\end{lemma}

\begin{proof}
\leanok
Divide the generators by $d_\nu$ so that they are coprime. Two consecutive elements then lie in
the monoid, and every sufficiently large integer is a nonnegative combination of them.
\end{proof}

\begin{definition}[Admissible descents]
\label{def:descend}
\lean{GraphMarkovMatching.Descend}
\leanok
\uses{def:grammar}
Write $t\Downarrow_m t'$ when there is a chain of $m$ successor steps from $t$ to $t'$.
\end{definition}

\begin{lemma}[Realisation and decomposition of fresh returns]
\label{thm:descend-realise}
\lean{GraphMarkovMatching.descend_of_mem_toFresh}
\leanok
\uses{def:descend, def:renewal}
The fresh-to-fresh descent lengths are exactly the monoid: $\set{m\colon\Ff\Downarrow_m\Ff}=\Sigma_\nu$. Every letter admits some descent to $\Ff$, and every return length in
$\mathrm{toF}(k)$ with $k\in S_\nu$ is realised.
\end{lemma}

\begin{proof}
\leanok
\uses{def:renewal}
For realisation, induct along the return-length recursion: whenever the successor set of a
letter contains the cascade pair of $k$, each $d\in\mathrm{toF}(k)$ is realised by a descent of that
length. For decomposition, every letter has some descent to $\Ff$, and concatenating fresh returns
shows the descent lengths form the monoid.
\end{proof}

\begin{lemma}[List coverage]
\label{thm:coverage}
\lean{GraphMarkovMatching.mem_zsucc_iter_of_descend}
\leanok
\uses{def:descend}
If $t\in\boldsymbol z$ and $t\Downarrow_mt'$, then $t'$ lies in the $m$-fold descended list, and
along any chain of full successor screens the zero list after $n$ steps is the $n$-fold descent
of the initial one.
\end{lemma}

\begin{proof}
\leanok
\uses{def:descend}
The empty descent leaves $t'=t$ in the list. One step places $t'$ in the descended list, and
induction advances the membership by one iterate per step; the same induction runs along a chain of
full successor screens.
\end{proof}

\begin{lemma}[Returns along source chains]
\label{thm:source-ray}
\lean{GraphMarkovMatching.hits_F_of_gpair}
\leanok
Along a chain of successor letters, a cascade entry at counter $k$ returns to $\Ff$ after some
$m\in\mathrm{toF}(k)$, and a fresh letter returns after some renewal depth.
\end{lemma}

\begin{proof}
\leanok
\uses{thm:descend-realise}
Strong induction on the counter. If the next letter is $\Ff$ then $1\in\mathrm{toF}(k)$ works;
otherwise the next letter is a forced letter at a strictly smaller counter, and the inductive
hypothesis supplies a return whose length lies in the corresponding shifted set.
\end{proof}

\begin{definition}[Anchored letters and screens]
\label{def:anchored}
\lean{GraphMarkovMatching.Anchored}
\leanok
\uses{def:descend, thm:semigroup}
A letter is anchored at absolute depth $d$ when $t\Downarrow_m\Ff$ implies $d+m\in\Sigma_\nu$. A
screen is anchored at $d$ when its cell and every member of its zero list are.
\end{definition}

\begin{lemma}[Anchoring propagates]
\label{thm:anchor-step}
\lean{GraphMarkovMatching.Anchored.step}
\leanok
\uses{def:anchored}
If $t$ is anchored at $d$ and $t'$ is a successor, then $t'$ is anchored at $d+1$; the same for
screens. If $\Ff$ is anchored at $d$ then $d\in\Sigma_\nu$, and every accessible pair has both
coordinates anchored.
\end{lemma}

\begin{proof}
\leanok
\uses{def:anchored}
Prepending the step $t\to t'$ to a fresh completion of $t'$ of length $m$ gives one of $t$ of
length $m+1$, which is the propagation. The empty completion forces $d\in\Sigma_\nu$ when $\Ff$ is
anchored at $d$, and the accessible pairs are anchored because the seed is.
\end{proof}

\begin{theorem}[No anchored live path]
\label{thm:nilpotence}
\lean{GraphMarkovMatching.no_anchored_live_path}
\leanok
\uses{def:scracc}
There is no infinite chain of formal screens, each a full successor of the last, starting from
an anchored screen with nonempty zero list and never having its cell in its zero list.
Consequently the block matrix on the live index is nilpotent.
\end{theorem}

\begin{proof}
\leanok
\uses{thm:anchor-step, thm:coverage, thm:source-ray}
Suppose such a chain exists. Anchoring propagates along it by \cref{thm:anchor-step}, and
\cref{thm:source-ray} returns the cell to $\Ff$ at a bounded number of steps. \Cref{thm:coverage}
puts that same letter in the zero list at the same step, so the cell lies in its zero list,
contradicting liveness. A finite live index with no infinite path makes the block matrix
nilpotent.
\end{proof}

\section{The closed recursion}

\begin{definition}[Ordinary scalar and screen vector]
\label{def:running}
\lean{GraphMarkovMatching.genPsi}
\leanok
\uses{def:ordacc, def:scracc}
Put $\overline\Psi_h:=\sup_{(s,t)\in\mathfrak O_\nu}\Psi_h(s,t)$, let $E_h$ be the vector of
screens over the live index, and $\overline E_h$ its supremum.
\end{definition}

\begin{definition}[Ordinary and screen step functions]
\label{def:step-functions}
\lean{GraphMarkovMatching.genF, GraphMarkovMatching.genG}
\leanok
\uses{def:cellCB, def:oneSided}
Fix a tilt bound and ledger bounds. The ordinary step function $f$ and screen step function $g$
are the explicit increasing functions assembled from the cell and one-sided bounds.
\end{definition}

\begin{lemma}[The ordinary step]
\label{thm:psi-step}
\lean{GraphMarkovMatching.genPsi_step}
\leanok
\uses{def:step-functions}
$\overline\Psi_{h+1}\leq f(\overline\Psi_h,\overline E_h)$ for every $h$.
\end{lemma}

\begin{proof}
\leanok
\uses{thm:psi-rows, thm:ordacc-inv}
An accessible pair contains no component letter, so its coordinate at height $h+1$ is one of the
four rows of \cref{thm:psi-rows}, and every row bound is dominated by the ordinary step
function.
\end{proof}

\begin{lemma}[The screen step]
\label{thm:e-step}
\lean{GraphMarkovMatching.genE_step}
\leanok
\uses{def:accN, def:step-functions}
For $\mu(0)>0$ and $C_W$ large enough,
$E_{h+1}(i)\leq g(\overline\Psi_h,\overline E_h)+(\Nb E_h)(i)$ for every $h$ and every $i$.
\end{lemma}

\begin{proof}
\leanok
\uses{thm:screen-rows, thm:scracc-inv}
An index screen has one of the six shapes of \cref{thm:screen-rows}, by the component-letter
freeness of \cref{thm:scracc-inv}. Each row bounds the screen at height $h+1$ by the step function
plus the sum over full successors, which is the block matrix applied to $E_h$.
\end{proof}

\begin{lemma}[Nilpotent invariant]
\label{thm:invariant}
\lean{GraphMarkovMatching.nilpotent_invariant}
\leanok
Let $\Nb$ be nonnegative on a finite index with $\Nb^{\,r}(u\mathbf1)=0$, and put
$\widehat E:=\sum_{j<r}\Nb^{\,j}(u\mathbf1)$. Then $u\mathbf1+\Nb\widehat E=\widehat E$, and any
vector sequence dominated by the step stays below $\widehat E$.
\end{lemma}

\begin{proof}
\leanok
\uses{thm:nilpotence}
The identity is the index shift, whose top term $\Nb^{\,r}(u\mathbf1)$ vanishes. Invariance
follows by induction on $h$: the base is the definition, and one step applies the monotone bound and
the identity.
\end{proof}

\begin{theorem}[The screened recursion closes]
\label{thm:engine}
\lean{GraphMarkovMatching.screened_uniform_bound_mono}
\leanok
Under nilpotence of $\Nb$, monotone step functions, and the size and closure inequalities, the
ordinary scalar stays below $K_c\eta$ and the screen vector below $\widehat E$ at every height.
\end{theorem}

\begin{proof}
\leanok
\uses{thm:invariant, thm:psi-step, thm:e-step}
Induct on the height on the joint invariant that the scalar stays below $K_c\eta$ and the vector
below $\widehat E$. The base is the size inequality, and the step applies the two monotone step
functions together with the invariant identity of \cref{thm:invariant}.
\end{proof}

\begin{lemma}[Geometric bound for the invariant]
\label{thm:geom}
\lean{GraphMarkovMatching.mulVec_le_of_row_sum}
\leanok
If every row sum of $\Nb$ is at most $c$, then $\Nb^{\,j}(u\mathbf1)\leq c^ju\mathbf1$ and the
invariant is at most $(\sum_{j<r}c^j)u\mathbf1$. For the block matrix one may take
$c=C_W\abs{\mathfrak I_\nu}$.
\end{lemma}

\begin{proof}
\leanok
\uses{thm:invariant}
If $x\leq b\mathbf1$ entrywise then $(\Nb x)(i)\leq b\sum_j\Nb(i,j)\leq cb$ for every $i$, so one
application multiplies a constant bound by at most $c$. Induction on $j$ gives the power bound and
summing gives the invariant bound.
\end{proof}

\begin{theorem}[Uniform failure bound, closure form]
\label{thm:closure-form}
\lean{GraphMarkovMatching.varying_failure_uniform}
\leanok
With the step functions of \cref{def:step-functions} and $r_\nu=\abs{\mathfrak I_\nu}+1$ from
\cref{thm:nilpotence}, and the ledger bounds in force, the failure probability is bounded
uniformly in the height by the closed form.
\end{theorem}

\begin{proof}
\leanok
\uses{thm:engine, thm:geom, thm:base, thm:nilpotence}
Apply \cref{thm:engine} to the running coordinates, with the block matrix of \cref{def:accN},
the step functions of \cref{def:step-functions}, nilpotence from \cref{thm:nilpotence}, the base
bounds of \cref{thm:base}, and the invariant bounded geometrically by \cref{thm:geom}.
\end{proof}

\section{The numerical closure}

\begin{lemma}[Ledger constants]
\label{thm:ledger-constants}
\lean{GraphMarkovMatching.rootTilt_le}
\leanok
If $\sim$ is symmetric and $\mu(0)\geq1/2$, the ledger hypotheses hold with $R_T=8$,
$F_M=2\eta$, $F_T=2\eta$.
\end{lemma}

\begin{proof}
\leanok
\uses{thm:zero-interface}
For $v\sim0$ the ball $B_G(v,1)$ contains $0$, so $b(v)\geq\mu(0)\geq\tfrac12$ and
$b(v)^{-5/2}\leq2^{5/2}\leq8$. The far mass and the far tilt are the two zero-interface bounds of
\cref{thm:zero-interface}.
\end{proof}

\begin{definition}[The explicit constants]
\label{def:constants}
\lean{GraphMarkovMatching.genKcC}
\leanok
\uses{thm:ledger-constants, thm:universes}
The constants $C_W$, $K_c$ and the level bounds are fixed from the tilt bound, the cardinality
bounds $c_N\geq\abs{\mathfrak I_\nu}$ and $n_{\cA}\geq\abs{\cA_N}$, and the support size.
\end{definition}

\begin{theorem}[Numerical closure]
\label{thm:numeric}
\lean{GraphMarkovMatching.varying_failure_le}
\leanok
\uses{def:constants}
If $\sim$ is reflexive and symmetric, $\mu(0)\geq1/2$ and $\eps^{-1}\eta\leq1$, then the size
and closure inequalities hold with the explicit constants and the failure probability is at most
$K_c\eta$.
\end{theorem}

\begin{proof}
\leanok
\uses{thm:closure-form, thm:ledger-constants}
Instantiate \cref{thm:closure-form} at the constants of \cref{def:constants} with the ledger
values of \cref{thm:ledger-constants}, and verify the size and closure inequalities in turn using
$\eps^{-1}\eta\leq1$.
\end{proof}

\begin{lemma}[Marginal consistency]
\label{thm:marginal-consistency}
\lean{GraphMarkovMatching.Tlaw_map_restrictLab}
\leanok
Restriction of a height-$(h+1)$ labelling to $\bbB_h$ pushes $T_{h+1}$ forward to $T_h$, so the
finite-height bounds pass to the infinite tree.
\end{lemma}

\begin{proof}
\leanok
The restriction of a branch is the branch of the restrictions, so the Markov property gives the
identity by induction on the height, one fibre of the root mixture at a time.
\end{proof}

\section{Two counter laws with a common core}

\begin{definition}[Composite kernels]
\label{def:composite-kernel}
\lean{GraphMarkovMatching.Composite.compK}
\leanok
\uses{def:grammar}
The label alphabet consists of the plain labels $(v,k)$ together with the marked labels
$(0,\langle p,i\rangle)$ for a declared pair $p$ and a stage $i$ along its designated path. The
composite kernel walks the designated path and rejoins the common core.
\end{definition}

\begin{lemma}[Cluster identity]
\label{thm:cluster}
\lean{GraphMarkovMatching.Composite.mul_repZ_le_cZ}
\leanok
\uses{def:composite-kernel}
For a side $\sigma$, a declared pair $p=(a,b)$ and a height $h$, the marked pair law is the
cascade law at $a$ conditioned on the designated port carrying the label $(0,b)$, an event of
probability $\mu(0)\nu_\sigma(b)>0$.
\end{lemma}

\begin{proof}
\leanok
\uses{def:composite-kernel}
Below a counter-$a$ root the balanced rule forces every label of the $a$-cascade, and the
designated port carries a fresh sample independent of all others. Conditioning on that sample being
$(0,b)$ is an event of probability $\mu(0)\nu_\sigma(b)$ and leaves the marked law.
\end{proof}

\begin{lemma}[Support equality]
\label{thm:support-equality}
\lean{GraphMarkovMatching.Composite.exists_sim_of_muM_ne_zero}
\leanok
\uses{def:composite-kernel}
After forgetting the counters, the supports of the corresponding components on the two sides
agree, componentwise and at every height.
\end{lemma}

\begin{proof}
\leanok
Induct on the height, proving all identities simultaneously. At height zero every listed law has
root-state law $\mu$ or the point mass at $0$, and the branch recursion preserves the componentwise
agreement because the two kernels differ only in the counters.
\end{proof}

\begin{corollary}[Exact pruning across the two laws]
\label{thm:exact-pruning}
\lean{GraphMarkovMatching.Composite.dead_screen_eq_zero}
\leanok
A screen whose cell is a component on one side and whose zero list holds the matching components
on the other side vanishes.
\end{corollary}

\begin{proof}
\leanok
\uses{thm:support-equality, thm:pruning}
By \cref{thm:support-equality} every point of positive source mass has a target point carrying
the same states at every vertex, and $\sim$ is reflexive, so the identity pairing witnesses life.
The dead-event indicator therefore vanishes, as in \cref{thm:pruning}.
\end{proof}

\begin{definition}[The tagged alphabet and letter successors]
\label{def:composite-grammar}
\lean{GraphMarkovMatching.Composite.Letter}
\leanok
\uses{def:grammar, def:composite-kernel}
The tagged alphabet adds to the ordinary letters the marker letters carrying a declared pair and
a stage, with successors following the designated path.
\end{definition}

\begin{definition}[The live index of the tagged grammar]
\label{def:composite-index}
\lean{GraphMarkovMatching.Composite.screenUniv}
\leanok
\uses{def:composite-grammar, def:scracc}
The seeds are the fresh screens in either orientation, and the live index is the live part of
their closure under full successors, restriction of the zero list, and the two side maps.
\end{definition}

\begin{lemma}[Universe bound]
\label{thm:composite-universe}
\lean{GraphMarkovMatching.Composite.cLetterBox_card}
\leanok
\uses{def:composite-index}
With $\cL_N$ the formal envelope of one fresh name, the counter names $0,\dots,N$ and all marker
names with coordinates and stage in $\set{0,\dots,N}$,
$\abs{\cL_N}=1+(N+1)+(N+1)^3=n_{\cA}$.
\end{lemma}

\begin{proof}
\leanok
\uses{thm:universes}
The three families of names are disjoint with the displayed cardinalities. A screen is a cell, a
subset of the alphabet and a normalisation, which gives $2^{n_{\cA}}$ zero lists and $n_{\cA}+1$
normalisations.
\end{proof}

\begin{lemma}[Chart returns]
\label{thm:chart-returns}
\lean{GraphMarkovMatching.Composite.chartPorts_cases}
\leanok
\uses{def:renewal, def:composite-kernel}
The designated port depth satisfies $d(k)\in\mathrm{toF}(k)$, and the port depths of the chart of
a declared pair lie in the corresponding union of return-length sets.
\end{lemma}

\begin{proof}
\leanok
\uses{def:renewal}
Induct along the port-depth recursion: $d(k)=1$ for $k\leq3$, and for $k\geq4$ the depth enters
the floor branch of the return-length recursion. The chart of a declared pair concatenates the
designated path with the exposed fresh ports, so its depths lie in the corresponding union.
\end{proof}

\begin{lemma}[Realisation]
\label{thm:composite-realise}
\lean{GraphMarkovMatching.Composite.descend_F_iff_mem_sigmaS}
\leanok
On either side the fresh-to-fresh descent lengths are $\Sigma_S$, the descent lengths to $\Ff$
do not depend on the side, and every chain of $2N+2$ successor steps passes through $\Ff$.
\end{lemma}

\begin{proof}
\leanok
\uses{thm:chart-returns, thm:descend-realise}
On either side the successor set of $\Ff$ contains the cascade pair of every core counter, so
the realisation half of \cref{thm:descend-realise} applies verbatim. \Cref{thm:chart-returns} places
the marker depths in the same monoid, so the descent lengths do not depend on the side.
\end{proof}

\begin{definition}[Anchored tagged letters and screens]
\label{def:composite-anchored}
\lean{GraphMarkovMatching.Composite.Anchored, GraphMarkovMatching.Composite.ScreenAnchored}
\leanok
\uses{def:anchored, thm:composite-realise}
A tagged letter is anchored at $\delta$ when $\delta+m\in\Sigma_S$ for every descent length $m$
to $\Ff$ on either side, and a tagged screen when its cell, zero list and normalisation are.
\end{definition}

\begin{lemma}[Anchoring propagates]
\label{thm:composite-anchor}
\lean{GraphMarkovMatching.Composite.Anchored.step}
\leanok
\uses{def:composite-anchored}
$\Ff$ is anchored at $\delta$ if and only if $\delta\in\Sigma_S$, and a successor of a letter
anchored at $\delta$ is anchored at $\delta+1$.
\end{lemma}

\begin{proof}
\leanok
\uses{thm:anchor-step, thm:composite-realise}
The two clauses are \cref{thm:anchor-step} with $\Sigma_S$ in place of $\Sigma_\nu$ and the
side-indexed successors in place of the ordinary ones; the empty chain forces $\delta\in\Sigma_S$
for the fresh letter.
\end{proof}

\begin{theorem}[Acyclicity of the screen steps]
\label{thm:composite-acyclic}
\lean{GraphMarkovMatching.Composite.no_live_step_cycleS}
\leanok
No cycle of full-successor edges passes through live anchored tagged screens, whether an edge
advances the cell by a common counter or through an exceptional entrance. The directed graph on
the live index is acyclic and the block matrix is nilpotent.
\end{theorem}

\begin{proof}
\leanok
\uses{thm:composite-anchor, thm:nilpotence}
A cycle through live anchored screens, followed forever, gives an infinite path of live anchored
screens, which \cref{thm:nilpotence} excludes once anchoring propagates by
\cref{thm:composite-anchor}. A finite acyclic index makes the block matrix nilpotent.
\end{proof}

\begin{lemma}[The weighted mixture tilt]
\label{thm:mixture-tilt-composite}
\lean{GraphMarkovMatching.Composite.invTilt_simDeg_cXiBar_le}
\leanok
\uses{def:composite-kernel}
At every pair point the restricted weight of the mixed pair law on one side is bounded by the
weighted sum over the common core and the exceptional entrances.
\end{lemma}

\begin{proof}
\leanok
\uses{thm:mixture-tilt}
Degrees are linear in the target law, so the mixed degree splits over the common core and the
exceptional entrances with the exact coefficients of the counter law. Dividing through, as in
\cref{thm:mixture-tilt}, gives the weighted bound.
\end{proof}

\begin{lemma}[Retention of the source coefficients]
\label{thm:retention}
Every transfer-row expression is a sum of integrals against its source law, so expanding the
fresh source cell splits the row into component rows with the exact coefficients of the counter
law, on both the common core and the exceptional entrances.
\end{lemma}

\begin{proof}
\uses{thm:mixture-tilt-composite}
Linearity of the source integral gives the split with the exact coefficients before any bound is
applied, and \cref{thm:mixture-tilt-composite} supplies the weight for each component row.
\end{proof}

\begin{proposition}[The composite ledger]
\label{thm:composite-ledger}
\lean{GraphMarkovMatching.Composite.cPsi_step}
\leanok
\uses{thm:retention, thm:composite-acyclic, thm:psi-step, thm:e-step}
For two counter laws with a common core, the running coordinates formed from the two orientations
of the composite comparison satisfy the ordinary and screen step inequalities of
\cref{thm:psi-step} and \cref{thm:e-step} with the tagged step functions.
\end{proposition}

\begin{proof}
\leanok
\uses{thm:retention, thm:psi-step, thm:e-step}
Each transfer row splits over the common core and the exceptional entrances with the exact
coefficients of the two counter laws, by \cref{thm:retention}. Applying the ordinary and screen
rows of \cref{thm:psi-step} and \cref{thm:e-step} to each component and recombining gives the
tagged step functions.
\end{proof}

\section{Proof of the two-law theorem}

\begin{proof}
\leanok
\proves{thm:composite-matching}
\uses{thm:composite-ledger, thm:composite-acyclic, thm:engine, thm:exact-pruning, thm:composite-universe}
The tagged grammar's block matrix is nilpotent by \cref{thm:composite-acyclic}, and the composite
ledger \cref{thm:composite-ledger} supplies the two step inequalities. \Cref{thm:engine} then
closes the recursion, and the cardinality bound of \cref{thm:composite-universe} fixes the
constants.
\end{proof}
